\documentclass[11pt,a4paper]{article}
\usepackage[margin=1in]{geometry}
\usepackage{amsmath,amssymb,amsthm}
\usepackage{microtype}
\usepackage[T1]{fontenc}
\usepackage{lmodern}
\usepackage[hidelinks]{hyperref}

\title{Homogeneous Ideals for the Named Linear Spaces\\ in $\mathbb{P}^5$}
\author{Benjamin Frohman}
\date{22 September 2026}

\begin{document}
\maketitle

\section{What can be written}

A Frohmanian \emph{label} does not create a new subvariety. The two objects named below are the coordinate plane and a complementary pair of planes in $\mathbb{P}^5$. Their ideals are the standard prime (resp.\ intersection) ideals of those linear spaces. Spacetime $\mathbb{R}^3\times\mathbb{R}$ is not a closed subvariety of $\mathbb{P}^5$ and has no homogeneous ideal in $\mathbb{C}[x_0,\ldots,x_5]$ until an embedding is chosen. None is chosen here.

\section{The plane}

Let
\[
\Pi=\bigl\{\,[0:0:0:x_3:x_4:x_5]\;\big|\; [x_3:x_4:x_5]\in\mathbb{P}^2\,\bigr\}
=\mathbb{V}(x_0,x_1,x_2)\subset\mathbb{P}^5.
\]
This is the statement $\Pi=\langle x_3,x_4,x_5\rangle$ in dual notation. The saturated homogeneous ideal is
\begin{equation}
\label{eq:IC}
I(\Pi)=(x_0,x_1,x_2)\subset\mathbb{C}[x_0,x_1,x_2,x_3,x_4,x_5].
\end{equation}
It is prime, generated by three linear forms, of height $3$ and of dimension $2$. The class $\gamma=[\Pi]$ is the fundamental class of this plane in $H_4(\mathbb{P}^5,\mathbb{Z})$, or, after restriction to a fourfold $X$ containing $\Pi$, the class of that surface in $H^4(X,\mathbb{Z})\cap H^{2,2}(X)$.

That ideal and that class are the coordinate plane in $\mathbb{P}^5$. They predate any personal name.

\section{The complementary pair on the split quadric}

Let
\[
\Pi'=\mathbb{V}(x_3,x_4,x_5).
\]
Then
\begin{equation}
\label{eq:ICp}
I(\Pi')=(x_3,x_4,x_5),
\end{equation}
and the union has
\begin{equation}
\label{eq:Iunion}
I(\Pi\cup\Pi')
=(x_0,x_1,x_2)\cap(x_3,x_4,x_5)
=(x_ix_j:i\in\{0,1,2\},\; j\in\{3,4,5\}).
\end{equation}
Both planes lie on the split quadric fourfold
\begin{equation}
\label{eq:Q}
Q^4=\mathbb{V}(x_0x_5+x_1x_4+x_2x_3)\subset\mathbb{P}^5,
\end{equation}
because each generator of $I(Q^4)$ vanishes on $\Pi$ and on $\Pi'$. Thus
\[
I(Q^4)=(x_0x_5+x_1x_4+x_2x_3).
\]
The pair $(\Pi,\Pi')$ is one pair of planes from the two rulings of this quadric. That configuration is classical.

\section{A sextic host containing $\Pi$}

A hypersurface $X\subset\mathbb{P}^5$ of degree $6$ contains $\Pi$ if and only if its equation lies in $I(\Pi)$. Equivalently,
\begin{equation}
\label{eq:sextic}
F=x_0A+x_1B+x_2C,\qquad A,B,C\in\mathbb{C}[x_0,\ldots,x_5]_5,
\end{equation}
and $X=\mathbb{V}(F)$. Then $I(X)=(F)$ (assuming $F$ irreducible) and $I(\Pi\subset X)$ is still~\eqref{eq:IC}, now viewed as an ideal of the homogeneous coordinate ring of $X$ after restriction. Smoothness of $X$ is an open condition on $(A,B,C)$ and is not automatic.

This writes a host containing the named plane. It does not select a single $X$ and it does not produce a new Hodge number.

\section{What is not written}

\begin{itemize}
\item There is no homogeneous ideal of spacetime $\mathbb{R}^3\times\mathbb{R}$ in $\mathbb{C}[x_0,\ldots,x_5]$. That four-manifold is affine, not a closed subvariety of $\mathbb{P}^5$.
\item There is no polynomial $F_{\mathrm{Frohman}}$ singled out by the Navier--Stokes work. The coadjoint orbit of $\mathrm{SDiff}$ and $\mathbb{T}^3\times[0,T_*)$ are not $\mathbb{P}^5$.
\item Writing~\eqref{eq:IC} or~\eqref{eq:Iunion} does not assign a novel Hodge number to $\eta_{F,l}$. The class $[\Pi]$ on a sextic fourfold, or $[\Pi]+[\Pi']$ on $Q^4$, has the intersection numbers of a linear plane (resp.\ of a pair of planes on a quadric). Those numbers are the ordinary ones.
\item A residue $\mathrm{Res}(\nu)$ and a Hodge number $h^{2,2}$ become available only after a single smooth $X$ is fixed and the primitive cohomology of that $X$ is computed. The family~\eqref{eq:sextic} is not a single $X$.
\end{itemize}

\section{Locked list, and only this}

\[
\begin{aligned}
I(\Pi)&=(x_0,x_1,x_2),\\
I(\Pi')&=(x_3,x_4,x_5),\\
I(\Pi\cup\Pi')&=(x_ix_j:i\in\{0,1,2\},\;j\in\{3,4,5\}),\\
I(Q^4)&=(x_0x_5+x_1x_4+x_2x_3).
\end{aligned}
\]
Until a single irreducible $F$ of degree $6$ is chosen, there is no $I(X_{\mathrm{sextic}})$ to lock, and $\eta_{F,l}$ remains a name without a numerical Hodge datum.

\end{document}
